% from ega3-1-fr.tex
\oldpage[III]{116}
l'égalité précédente résulte de (I, 5.2.1 et 1.1.5)~; si $z\notin Y''$, on a
donc $\mathcal J'_z=\mathcal J_z$, d'où $\mathcal F'_z=\mathcal F_z$ et
$\mathcal F''_z=0$, ce qui établit notre assertion dans ce cas~; on raisonne
de même si $z\notin Y'$. Par suite, le noyau et le conoyau de $u$, qui sont
dans $K$ (0\textsubscript{I}, 5.3.4) ont leurs supports dans $Y'\cap Y''$,
et sont donc dans $K'$ par hypothèse~; pour la même raison, $\mathcal F'$ et
$\mathcal F''$ sont dans $K'$, donc aussi $\mathcal F'\oplus\mathcal F''$,
puisque $K'$ est exact. La conclusion résulte alors de la considération des
deux suites exactes
\[
0\longrightarrow\operatorname{Im}u\longrightarrow
\mathcal F'\oplus\mathcal F''\longrightarrow
\operatorname{Coker}u\longrightarrow0,
\]
\[
0\longrightarrow\operatorname{Ker}u\longrightarrow\mathcal F
\longrightarrow\operatorname{Im}u\longrightarrow0
\]
et de l'hypothèse que $K'$ est exact.

\item[b)] $Y$ est irréductible, et par suite le sous-préschéma $Y$ de $X$ est
\emph{intègre}. Si $y$ est son point générique, on a
$(\mathcal O_Y)_y=k(y)$, et comme $j^*(\mathcal F)$ est un
$\mathcal O_Y$-Module cohérent, $\mathcal F_y=(j^*(\mathcal F))_y$ est un
$k(y)$-espace vectoriel de dimension finie $m$. Par hypothèse, il y a un
$\mathcal O_X$-Module cohérent $\mathcal G$ (nécessairement de support $Y$)
tel que $\mathcal G_y$ soit un $k(y)$-espace vectoriel de dimension $1$. Par
suite, il y a un $k(y)$-isomorphisme
$(\mathcal G_y)^m\xrightarrow{\sim}\mathcal F_y$, qui est aussi un
$\mathcal O_Y$-isomorphisme, et comme $\mathcal G^m$ et $\mathcal F$ sont
cohérents, il existe un voisinage ouvert $W$ de $y$ dans $X$ et un
isomorphisme $\mathcal G^m|W\xrightarrow{\sim}\mathcal F|W$
(0\textsubscript{I}, 5.2.7). Soit $\mathcal H$ le graphe de cet isomorphisme,
qui est un sous-$(\mathcal O_X|W)$-Module cohérent de
$(\mathcal G^m\oplus\mathcal F)|W$, canoniquement isomorphe à
$\mathcal G^m|W$ et à $\mathcal F|W$~; il existe donc un sous-$\mathcal O_X$-
Module cohérent $\mathcal H_0$ de $\mathcal G^m\oplus\mathcal F$, induisant
$\mathcal H$ sur $W$ et $0$ sur $X-Y$ puisque $\mathcal G^m$ et
$\mathcal F$ ont pour support $Y$ (I, 9.4.7). Les restrictions
$v:\mathcal H_0\to\mathcal G^m$ et $w:\mathcal H_0\to\mathcal F$ des
projections canoniques de $\mathcal G^m\oplus\mathcal F$ sont alors des
homomorphismes de $\mathcal O_X$-Modules cohérents, qui, dans $W$ et dans
$X-Y$, se réduisent à des isomorphismes~; autrement dit, les noyaux et
conoyaux de $v$ et $w$ ont leur support dans l'ensemble fermé
$Y-(Y\cap W)$, distinct de $Y$. Ils appartiennent donc à $K'$~; d'autre
part, on a $\mathcal G^m\in K'$ puisque $\mathcal G\in K'$ et que $K'$ est
exact. On en conclut successivement, par l'exactitude de $K'$, que
$\mathcal H_0\in K'$ puis que $\mathcal F\in K'$. C.Q.F.D.
\end{enumerate}
\end{proof}

\begin{corollary}[3.1.3]
\label{III.3.1.3-fr}
Supposons que le sous-ensemble exact $K'$ de $K$ ait en outre la propriété
que tout facteur direct cohérent d'un $\mathcal O_X$-Module cohérent
$\mathcal M\in K'$ appartienne encore à $K'$. Dans ces conditions, la
conclusion de (\hyperref[III.3.1.2-fr]{3.1.2}) subsiste encore lorsque la
condition «~$\mathcal G_y$ est un $k(y)$-espace vectoriel de dimension $1$~»
est remplacée par $\mathcal G_y\neq0$ (ce qui équivaut à
$\operatorname{Supp}(\mathcal G)=Y$).
\end{corollary}

\begin{proof}
En effet, le raisonnement de (\hyperref[III.3.1.2-fr]{3.1.2}) ne doit être
modifié que dans le cas b)~; cette fois $\mathcal G_y$ est un $k(y)$-espace
vectoriel de dimension $q>0$, et on a par suite un $\mathcal O_Y$-
isomorphisme $(\mathcal G_y)^m\xrightarrow{\sim}(\mathcal F_y)^q$~; la fin
du raisonnement de (\hyperref[III.3.1.2-fr]{3.1.2}) prouve alors que
$\mathcal F^q\in K'$, et l'hypothèse additionnelle sur $K'$ implique que
$\mathcal F\in K'$.
\end{proof}

\subsection{Le théorème de finitude : cas des schémas usuels}
\label{subsection:III.3.2-fr}

\begin{theorem}[3.2.1]
\label{III.3.2.1-fr}
Soient $Y$ un préschéma localement noethérien, $f:X\to Y$ un morphisme
propre. Pour tout $\mathcal O_X$-Module cohérent $\mathcal F$, les
$\mathcal O_Y$-Modules $R^qf_*(\mathcal F)$ sont cohérents pour $q\geq0$.
\end{theorem}

\begin{proof}
La question étant locale sur $Y$, on peut supposer $Y$ noethérien, donc $X$
noethérien (I, 6.3.7). Les $\mathcal O_X$-Modules cohérents $\mathcal F$ pour
lesquels la conclusion du th. (\hyperref[III.3.2.1-fr]{3.2.1}) est vraie
forment un sous-ensemble exact $K'$ de la catégorie $K$ des
$\mathcal O_X$-Modules cohérents.

